On prépare une solution aqueuse $(S)$ d'acide butanoïque \ce{C3H7COOH} de
concentration molaire $C=\qty{2,0e-3}{\M}$ et de volume $V=\qty{1,0}{\liter}$.
La mesure du pH de la solution $(S)$ a donné la valeur
$\mathrm{pH}=3,76$.
\begin{enumerate}
    \item Donner la définition d'un acide selon Brönsted.~\textbf{(1,0pt)}
    \item Écrire l'équation chimique modélisant la réaction de l'acide
          butanoïque avec l'eau.~\textbf{(1,0pt)}
    \item Déterminer la quantité de matière initiale de l'acide
          butanoïque~?~\textbf{(1,0pt)}
    \item Dresser le tableau d'avancement de la réaction.~\textbf{(1,0pt)}
    \item Déterminer la valeur de l'avancement maximal
          $x_{\max}$.~\textbf{(1,0pt)}
    \item Calculer la valeur de l'avancement à l'état final
          $x_{f}$.~\textbf{(1,0pt)}
    \item Calculer la valeur du taux d'avancement final $\tau$. Que peut-on
          déduire~?~\textbf{(1,0pt)}
\end{enumerate}

\begin{donnees}
    \begin{itemize}
    \item Couple acide/base~: \ch{C3H7COOH / C3H7COO-}.
\end{itemize}
\end{donnees}
